\section{A fixed finite quotient for the exact ordinary fibre}
\label{sec:G84}

This section sharpens the analytic specialization step in AC-G78 and
AC-G83. Throughout, \(p=269\), \(\mathcal O\) is the ring of integers
of a finite extension of \(\mathbf Q_p\), \(\pi\) is a uniformizer,
and \(B=\mathcal O[[W]]\). Absolute values satisfy \(|p|=p^{-1}\).
The actual ordinary branch and all finite Galois conditions of those
results remain in force.

\begin{proposition}[AC-G84-norm: fixed-precision root on a finite branch]
\label{prop:G84-norm}
Let \(R\) be a domain which is a finite torsion-free \(B\)-algebra,
and put \(u=1+p\). Define
\[
 f(W)=-\frac{(W-p)(W-(u^{-1}-1))}{1+W}.
 \tag{G84.1}
\]
Let \(S\) denote the integral closure of \(B\) in
\(\operatorname{Frac}(R)\). For every \(b\in R\) satisfying
\[
 b-f(W)\in(p^4,W^4)R
 \tag{G84.2}
\]
there are a finite extension \(E'/\operatorname{Frac}(\mathcal O)\)
and a continuous specialization
\[
 x:R\longrightarrow\mathcal O_{E'},\qquad
 x(b)=0,\qquad |x(W)-p|<p^{-2}.
 \tag{G84.3}
\]
More generally, the same conclusion holds under the weaker condition
\(b-f(W)\in(p^4,W^4)S\). The restriction of \(x\) to
\(\mathcal O\) is its given embedding.
The precision ideal is independent of the rank of \(R\), the
ramification of its weight map, and the ramification index of
\(\mathcal O/\mathbf Z_p\).
\end{proposition}

\begin{proof}
Let \(S\) be the integral closure of \(B\) in
\(\operatorname{Frac}(R)\). This fraction-field extension is finite.
The complete Noetherian local domain \(B\) is Nagata, so \(S\)
is finite over \(B\). It contains \(R\), is normal, and is
torsion-free over \(B\).
Every maximal ideal of \(S\) lies over \((\pi,W)\).
Going down over the normal domain \(B\), together with integrality,
shows that its localization has dimension two.
Normality gives Serre's condition \(S_2\), so these localizations
are Cohen--Macaulay. Since \(S/(\pi,W)S\) is finite-dimensional
over the residue field, \(\pi,W\) is a regular sequence on every
such localization, hence on \(S\).

Here is an explicit freeness argument. The Koszul resolution for
\(\pi,W\) over \(B\), after tensoring with \(S\), is exact in
positive degrees because this is an \(S\)-regular sequence. Thus
\(\operatorname{Tor}_1^B(B/(\pi,W),S)=0\).
Choose a finite free surjection \(B^d\to S\) inducing an isomorphism
modulo \((\pi,W)\), and let \(K\) be its kernel. The Tor sequence
gives \(K/(\pi,W)K=0\); Nakayama gives \(K=0\).
Consequently \(S\) is a finite free \(B\)-algebra of positive rank
\(d\). No etale point of the weight map has been assumed.

Set \(A=\mathcal O\langle Z\rangle\), the restricted power-series
ring, with the \(B\)-algebra structure
\[
 W=p+p^2Z.
 \tag{G84.4}
\]
The substitution converges in the \(\pi\)-adic topology.
The algebra \(S_A=S\otimes_B A\) is free of rank \(d\).
A direct computation gives
\[
 f(p+p^2Z)=p^3ZU(Z),\qquad
 U(Z)=-\frac{1+u^{-1}+pZ}{u+p^2Z}.
 \tag{G84.5}
\]
The series \(U\) is a unit of \(A\), reducing modulo \(\pi\)
to the nonzero constant \(-2\).
Also \(W^4=p^4(1+pZ)^4\) after substitution.
Hence (G84.2), or its stated weaker variant in \(S\), gives
\[
 b_A=p^3ZU+p^4a\quad\text{for some }a\in S_A.
 \tag{G84.6}
\]
Let \(M_a\in M_d(A)\) denote multiplication by \(a\). Then
\[
 \det(m_{b_A})=p^{3d}H(Z),\qquad
 H(Z)=\det(ZU\,I_d+pM_a),
 \tag{G84.7}
\]
and expansion by degree in \(p\) yields
\[
 H(Z)=Z^dU^d+
 \sum_{i=1}^d p^iZ^{d-i}U^{d-i}c_i(Z),\qquad c_i\in A.
 \tag{G84.8}
\]
Thus \(H/U^d\equiv Z^d\pmod{\pi A}\). No inverse trace matrix
or discriminant bound occurs in this calculation.

For completeness, the restricted-series factorization used here
can be constructed directly. If \(J\in A\) reduces to \(Z^d\),
construct monic polynomials \(P_n\) of degree \(d\), with lower
coefficients in \(\pi\mathcal O\), and \(V_n\in1+\pi A\), such that
\(J-P_nV_n\in\pi^nA\). Start with \(P_1=Z^d,V_1=1\).
The next error divided by \(\pi^n\), reduced modulo \(\pi\), is
a polynomial. Write it uniquely as
\(\overline L+Z^d\overline Q\) with \(\deg L<d\), and lift \(L,Q\)
to polynomials over \(\mathcal O\).
Replacing \(P_n,V_n\) by \(P_n+\pi^nL,V_n+\pi^nQ\) removes the
error modulo \(\pi^{n+1}\), since \(P_n\equiv Z^d\) and
\(V_n\equiv1\pmod\pi\). Completeness gives \(J=PV\),
where \(P\) is monic of degree \(d\), \(P\equiv Z^d\pmod\pi\),
and \(V\in1+\pi A\) is a unit.
Apply this to \(J=H/U^d\).
Every root \(z\) of \(P\) belongs to a finite coefficient extension
and satisfies \(|z|<1\): monicity and integrality give \(|z|\leq1\),
while \(|z|=1\) contradicts the strict dominance of the leading
term over all lower terms.

Choose a root \(z\), put
\(E_1=\operatorname{Frac}(\mathcal O)(z)\), and form
\[
 C=S_A\otimes_{A,Z\mapsto z}E_1.
\]
This is a nonzero algebra of dimension \(d\) over \(E_1\).
Multiplication by the image of \(b\) has determinant zero, so that
element is a nonunit and lies in a maximal ideal of \(C\).
Its residue field is a finite extension \(E'/E_1\).
The resulting map \(S\to E'\) kills \(b\).
Every image is integral over the image of \(B\), which lies in
\(\mathcal O_{E'}\); hence it lies in \(\mathcal O_{E'}\).
The map is continuous for the finite \(B\)-module topology:
\(\pi,W\) have topologically nilpotent images and the images of
finitely many module generators are bounded.
Its restriction to \(R\) proves (G84.3).
This last step produces an actual branch point even when the
norm root or the weight fibre is nonreduced.
\end{proof}

\begin{theorem}[AC-G84: one fixed finite trace quotient suffices]
\label{thm:G84}
Retain all actual ordinary-family, scalar-element, host,
projector and finite Galois hypotheses proved and used in
AC-G78--AC-G83, in the respective applicable branch.
Let \(R\) be their finite torsion-free ordinary branch, on the
fixed Teichmuller component and with coordinate \(W\) of AC-G78.
Write \(T:G_{\mathbf Q}\to R\) for its normalized rank-two trace.
Use the convention
\(\epsilon(\operatorname{Frob}_q)=q^{-1}\).
The actual scalar-corrected element \(g\) constructed there satisfies
\[
 \epsilon(g)=-(1+p),\qquad
 T(g)=-\big((1+W)+(1+W)^{-1}\big).
 \tag{G84.9}
\]
The exact full-chain ordinary fibre can be obtained by requiring
only the following additional trace precision:
\[
 T(\operatorname{Frob}_q)\equiv T(g)
       \pmod{(p^4,W^4)R},\qquad
 \epsilon(\operatorname{Frob}_q)\equiv\epsilon(g)\pmod{p^4}.
 \tag{G84.10}
\]
In the first congruence it is enough to use
\((p^4,W^4)S\), where \(S\) is the normalization just defined.
No equality \((p^4,W^4)S\cap R=(p^4,W^4)R\) is asserted.
These conditions can be imposed simultaneously with the actual
finite Kummer quotient and all the host, character and local
conditions of AC-G83. They yield a finite-coefficient point of
the branch with
\[
 T(\operatorname{Frob}_q)=q+q^{-1},\qquad
 |W-p|<p^{-2}.
 \tag{G84.11}
\]
The finite-Hecke-image argument of AC-G78 puts this point in the
actual ordinary module, with every from-the-start projector
retained. Its normalized three-dimensional Frobenius polynomial
is \((X-1)(X-q)(X-q^2)\).
\end{theorem}

\begin{proof}
The ring \(R/(p^4,W^4)R\) is finite because \(R\) is finite over
\(B\) and \(B/(p^4,W^4)\) is finite.
A continuous map from a profinite group to a finite discrete
set factors through some finite quotient by an open normal
subgroup. Indeed each fibre is compact and open, hence a finite
union of cosets of open subgroups; intersecting their finitely
many normal cores proves the assertion.
Apply this either to the reduction of \(T\) in \(R/(p^4,W^4)R\)
or to its composite in \(S/(p^4,W^4)S\); the latter ring is finite
because \(S\) is finite over \(B\). Intersect the resulting
subgroup with the kernels of \(\epsilon\bmod p^4\), the finite
host and character maps, and the fixed Kummer extension in AC-G83.
The conjugacy class of the actual element \(g\) in this single
finite Galois quotient therefore prescribes all conditions.

Chebotarev supplies primes outside the fixed finite exceptional
set in that class. The trace is conjugacy-invariant, and the
remaining conditions are exactly the compatible finite
conditions already checked in AC-G83. They retain splitting in
the designated host, the required cubic value,
\(p\parallel q+1\), and the unit-power condition for the real
dihedral cover. No independent, possibly conflicting Frobenius
prescription is added.

For such a prime set
\[
 b_q=T(\operatorname{Frob}_q)-
 \big(\epsilon(\operatorname{Frob}_q)
       +\epsilon(\operatorname{Frob}_q)^{-1}\big)\in R.
 \tag{G84.12}
\]
Inversion preserves congruences modulo \(p^4\) between units.
Writing \(u=1+p\), (G84.9) gives the scalar target
\[
 -\big((1+W)+(1+W)^{-1}\big)+(u+u^{-1})
   =-\frac{(W-p)(W-(u^{-1}-1))}{1+W}.
 \tag{G84.13}
\]
Thus \(b_q-f(W)\in(p^4,W^4)R\), or in
\((p^4,W^4)S\) for the normalization variant.
Proposition~\ref{prop:G84-norm} supplies an actual specialization
which kills \(b_q\), proving (G84.11). In the normalization
variant this is first an \(S\)-point and is then restricted to \(R\).
At this split prime the rank-two determinant is one, so the roots
are \(q,q^{-1}\). Adding the middle character one and using the
exact \(\epsilon^{-1}\) normalization from AC-G78 gives
\(q^2,q,1\).

This point belongs to the same finite branch \(R\), not merely to
weight space. The already proved map from the finite Hecke image
to \(R\), its arbitrary-base-change Cayley--Hamilton argument,
and the actual ordinary specialization theorem of AC-G78 apply.
All tame and coefficient projectors are present before
specialization, as in AC-G79 and AC-G80.
Hence the exact fibre is in the required projected ordinary
module. The subsequent four-test, projective finite-type and
real-cover arguments need no further analytic precision.
\end{proof}

\begin{remark}[Exact scope and source interface]
\label{rem:G84-scope}
The fixed ideal \((p^4,W^4)\), together with cyclotomic precision
\(p^4\), replaces an unspecified sufficiently fine approximation
near a possibly ramified weight-map point.
This does not bound the degree or discriminant of the finite
Galois quotient through which the reduced trace factors.
It supplies no numerical bound for \(q\), no counting-budget
estimate, and no fixed-level containment of an entire
characteristic-\(p\) deformation curve in a Hecke support.

The general algebra inputs are finite normalization for a
complete Noetherian domain, normality implying \(S_2\), and the
regular-sequence criterion in a Cohen--Macaulay local ring.
Exact primary locators are Stacks Project Lemma 10.162.8
(Tag 032W), Proposition 10.162.16 (Tag 0335), Proposition 10.38.7 (Tag 00H8), Lemma 10.157.4
(Tag 031S), and Lemma 10.104.2 (Tag 02JN), whose applied proof
is Proposition 10.103.4 (Tag 00N6).
The Koszul--Nakayama argument and restricted-series factorization
are proved above. The ordinary and automorphic interfaces are
exactly those audited in AC-G78--AC-G83; no new control or
automorphy theorem is asserted.
\end{remark}
